\documentclass[a4paper,12pt]{article}
\usepackage[T2A]{fontenc}
\usepackage[utf8]{inputenc}
\usepackage[russian]{babel}
\usepackage{amsmath,amssymb,mathtools}
\usepackage{helvet}
\renewcommand{\familydefault}{\sfdefault}
\usepackage{icomma}
\usepackage[margin=21mm,top=23mm,bottom=23mm,headheight=15pt]{geometry}
\usepackage{microtype}
\usepackage{tikz}
\usetikzlibrary{arrows.meta,positioning,shapes.geometric,calc}
\usepackage{tkz-euclide}
\usepackage{tabularx,booktabs}
\usepackage{enumitem}
\usepackage{fancyhdr}
\usepackage{setspace}
\usepackage{xcolor}
\usepackage{hyperref}
\definecolor{accent}{HTML}{008080}
\definecolor{darkaccent}{HTML}{004D4D}
\definecolor{textgray}{HTML}{343434}
\definecolor{muted}{HTML}{646464}
\hypersetup{hidelinks,pdftitle={Чек-лист: основные формулы тригонометрии},pdfauthor={Лёвин Артём Александрович}}
\newcommand{\keyidea}[1]{\par\smallskip\noindent\textcolor{darkaccent}{\textbf{Ключевая идея.}} #1\par\smallskip}
\newcommand{\alertline}[1]{\par\smallskip\noindent\textcolor{darkaccent}{\textbf{Проверьте.}} #1\par\smallskip}
\newcommand{\sectionheading}[1]{\section*{\textcolor{darkaccent}{#1}}}
\setlength{\parindent}{0pt}
\setlength{\parskip}{5pt}
\onehalfspacing
\setlist[itemize]{leftmargin=7mm,itemsep=2pt,topsep=2pt}
\setlist[enumerate]{leftmargin=9mm,itemsep=5pt,topsep=3pt}
\pagestyle{fancy}
\fancyhf{}
\fancyhead[L]{\footnotesize\textcolor{muted}{Основные формулы тригонометрии}}
\fancyhead[R]{\footnotesize\textcolor{muted}{Денис Ермолаев}}
\fancyfoot[C]{\footnotesize\textcolor{muted}{\thepage}}
\renewcommand{\headrulewidth}{0.25pt}
\renewcommand{\footrulewidth}{0pt}

\begin{document}
\begin{titlepage}
\thispagestyle{empty}
\begin{center}
\vspace*{24mm}
{\large\textcolor{darkaccent}{МАТЕМАТИКА}}\par
\vspace{24mm}
{\Huge\mdseries Чек-лист}\par
\vspace{13mm}
{\LARGE\textcolor{accent}{Основные формулы\\[2mm] тригонометрии}}\par
\vspace{13mm}
{\normalsize Подготовка к ЕГЭ по профильной математике}\par
\vspace{36mm}
{\large Денис Ермолаев}\par
\vspace{13mm}
{\large 16 ноября 2026 г.}\par
\vfill
{\small Лёвин Артём Александрович}\par
{\small эксклюзивно для Дениса Ермолаева}\par
\vspace{21mm}
\begin{tikzpicture}
\draw[accent,line width=0.85pt] (-3.2,0) .. controls (-2.0,0.62) and (-1.15,-0.45) .. (0,0)
 .. controls (1.15,0.50) and (2.0,-0.52) .. (3.2,0);
\end{tikzpicture}
\end{center}
\end{titlepage}

\sectionheading{1. Основное тождество и связанные функции}
Главное тождество справедливо для любого действительного угла $\alpha$:
\[
\boxed{\sin^2\alpha+\cos^2\alpha=1.}
\]
Оно позволяет восстановить модуль синуса по косинусу и наоборот:
\[
\cos\alpha=\pm\sqrt{1-\sin^2\alpha},\qquad
\sin\alpha=\pm\sqrt{1-\cos^2\alpha}.
\]
\keyidea{Знак корня выбирается по четверти, в которой расположен угол. Знаки \(\pm\) означают два возможных значения до учёта четверти.}
Любое число $c$ можно при необходимости переписать как
\[
 c=c(\sin^2\alpha+\cos^2\alpha)=c\sin^2\alpha+c\cos^2\alpha.
\]
\alertline{В тождестве у синуса и косинуса должен стоять \textbf{один и тот же угол}.}
Тангенс и котангенс связаны с синусом и косинусом:
\[
\tg\alpha=\frac{\sin\alpha}{\cos\alpha}\quad(\cos\alpha\ne0),
\qquad
\ctg\alpha=\frac{\cos\alpha}{\sin\alpha}\quad(\sin\alpha\ne0).
\]
\textbf{Пример.} Пусть $\sin\alpha=\dfrac35$ и $\dfrac\pi2<\alpha<\pi$. Тогда угол во второй четверти,
поэтому косинус отрицателен:
\[
\cos\alpha=-\sqrt{1-\frac9{25}}=-\frac45,\qquad
\tg\alpha=-\frac34,\quad\ctg\alpha=-\frac43.
\]

\sectionheading{2. Четверти, знаки и тригонометрическая окружность}
Положительное направление обхода окружности --- \textbf{против часовой стрелки};
отрицательное --- \textbf{по часовой стрелке}.
\begin{center}
\begin{minipage}[c]{0.39\linewidth}
\centering
\begin{tikzpicture}[scale=1.25,>=Stealth]
\tkzDefPoint(0,0){O}
\tkzDefPoint(1.4,0){E}
\tkzDrawCircle[line width=.75pt,color=darkaccent](O,E)
\draw[->,thin](-1.75,0)--(1.78,0)node[right]{$x$};
\draw[->,thin](0,-1.65)--(0,1.75)node[above]{$y$};
\node[font=\scriptsize] at (0.82,0.83) {I};
\node[font=\scriptsize] at (-0.84,0.83) {II};
\node[font=\scriptsize] at (-0.84,-0.83) {III};
\node[font=\scriptsize] at (0.82,-0.83) {IV};
\draw[->,accent,line width=1pt] (0.55,0) arc[start angle=0,end angle=72,radius=.55];
\node[font=\scriptsize,accent] at (1.35,.50) {$+\alpha$};
\end{tikzpicture}
\end{minipage}\hfill
\begin{minipage}[c]{0.57\linewidth}
\centering
\renewcommand{\arraystretch}{1.24}
\begin{tabular}{@{}lcccc@{}}
\toprule
Функция & I & II & III & IV\\
\midrule
$\sin$ & $+$ & $+$ & $-$ & $-$\\
$\cos$ & $+$ & $-$ & $-$ & $+$\\
$\tg,\ \ctg$ & $+$ & $-$ & $+$ & $-$\\
\bottomrule
\end{tabular}
\end{minipage}
\end{center}
\textbf{Границы четвертей:} на осях знак функции может быть нулевым,
а тангенс или котангенс может оказаться не определён.
\keyidea{Если задан интервал для угла, сначала установите четверть, затем выберите знак вычисленного значения.}

\sectionheading{3. Синус и косинус суммы и разности}
Для любых действительных $a,b$:
\begin{align*}
\sin(a+b)&=\sin a\cos b+\cos a\sin b,\\
\sin(a-b)&=\sin a\cos b-\cos a\sin b,\\
\cos(a+b)&=\cos a\cos b-\sin a\sin b,\\
\cos(a-b)&=\cos a\cos b+\sin a\sin b.
\end{align*}
\textbf{Пример.} Вычислим $\sin75^\circ$, разбив угол на два табличных:
\begin{align*}
\sin75^\circ&=\sin(45^\circ+30^\circ)\\
&=\sin45^\circ\cos30^\circ+\cos45^\circ\sin30^\circ\\
&=\frac{\sqrt2}{2}\cdot\frac{\sqrt3}{2}+\frac{\sqrt2}{2}\cdot\frac12
=\frac{\sqrt6+\sqrt2}{4}.
\end{align*}
\alertline{Для синуса знак между слагаемыми сохраняется; для косинуса в формуле суммы появляется минус.}

\sectionheading{4. Формулы двойного угла}
\[
\boxed{\sin2a=2\sin a\cos a}
\]
\[
\boxed{\cos2a=\cos^2a-\sin^2a
=2\cos^2a-1=1-2\sin^2a.}
\]
\textbf{Работа в обе стороны.} Если нужно свернуть произведение, ищите два одинаковых угла:
\[
\sin3a\cos3a=\frac12\sin6a.
\]
Если множитель $2$ отсутствует, его компенсируют делением на $2$.
\textbf{Пример.} При $\cos3a\ne0$:
\[
\frac{\sin6a}{\cos3a}
=\frac{2\sin3a\cos3a}{\cos3a}=2\sin3a.
\]

\sectionheading{5. Формулы приведения}
Большой угол часто удобно представить как кратный $\pi$ или $\pi/2$ плюс
или минус острый угол $\beta$.
\begin{enumerate}
\item Найдите конечную четверть на окружности с учётом \textbf{знака каждого слагаемого}.
\item Если базовый угол кратен $\pi$, \textbf{название функции сохраняется}.
Если он равен нечётному кратному $\pi/2$, функция меняется на кофункцию:
$\sin\leftrightarrow\cos$, $\tg\leftrightarrow\ctg$.
\item Знак результата определяется знаком \textbf{исходной функции} в конечной четверти.
\end{enumerate}
\textbf{Полезные преобразования} (при определённости обеих частей):
\begin{align*}
\sin(2\pi+\beta)&=\sin\beta,&
\cos(\pi+\beta)&=-\cos\beta,\\
\sin(\tfrac{3\pi}{2}+\beta)&=-\cos\beta,&
\tg(\tfrac\pi2+\beta)&=-\ctg\beta.
\end{align*}
\textbf{Пример 1.} $\sin(-10\pi+\pi/4)=\sin(\pi/4)=\sqrt2/2$.
Здесь $-10\pi$ означает пять полных оборотов по часовой стрелке,
а $+\pi/4$ --- поворот против неё.

\textbf{Пример 2.} $-\dfrac{7\pi}{2}+\dfrac\pi3
\equiv\dfrac{5\pi}{6}\pmod{2\pi}$, поэтому
\[
\cos\left(-\frac{7\pi}{2}+\frac\pi3\right)
=\cos\frac{5\pi}{6}=-\frac{\sqrt3}{2}.
\]
\alertline{В записи $-10\pi+\pi/4$ минус относится только к $10\pi$.
Для выражения $-(10\pi+\pi/4)$ результат поворота другой.}
\textbf{Область определения.} Равенство
$\tg(\pi/2+3a)=-\ctg3a$
используют при $\sin3a\ne0$.

\sectionheading{6. Частые ошибки и самопроверка}
\begin{itemize}[itemsep=0pt,topsep=0pt]
\item После извлечения корня оставлен знак $\pm$ без учёта четверти.
\item В тождестве $\sin^2a+\cos^2a=1$ перепутаны аргументы;
в формуле $\sin2a=2\sin a\cos a$ потерян множитель $2$.
\item При повороте знак первого слагаемого ошибочно перенесён на второе;
при сдвиге на нечётное $\pi/2$ забыта замена на кофункцию.
\item Сокращён знаменатель, который может обращаться в нуль.
\end{itemize}
\keyidea{Перед ответом проверьте четверть, знак, аргументы, множитель $2$ и ОДЗ.}

\clearpage
\thispagestyle{fancy}
\begin{center}
{\Large\textcolor{darkaccent}{Алгоритм: применение формул приведения}}\par
\vspace{11mm}
\begin{tikzpicture}[
 node distance=10mm and 16mm,
 >={Stealth[length=2.5mm]},
 stage/.style={draw=darkaccent,rounded corners=2mm,align=center,text width=10.1cm,minimum height=1.2cm,inner sep=5pt,font=\small},
 choice/.style={draw=accent,diamond,aspect=2.25,align=center,text width=4.6cm,inner sep=2pt,font=\small},
 branch/.style={draw=darkaccent,rounded corners=2mm,align=center,text width=4.7cm,minimum height=1.35cm,inner sep=5pt,font=\small},
 terminus/.style={draw=accent,ellipse,align=center,text width=6.6cm,minimum height=1.1cm,inner sep=3pt,font=\small},
 flow/.style={->,line width=0.8pt,draw=accent}
]
\node[terminus](begin){Дано: $F(k\pi/2\pm\beta)$, $\beta$ --- острый угол};
\node[stage,below=of begin](normalize){Упростите поворот на целое число полных оборотов $2\pi$};
\node[choice,below=of normalize](switch){Базовый угол --- нечётный кратный $\pi/2$?};
\node[branch,below left=13mm and 3mm of switch](yes){\textbf{Да:} замените функцию кофункцией\\$\sin\leftrightarrow\cos$, $\tg\leftrightarrow\ctg$};
\node[branch,below right=13mm and 3mm of switch](no){\textbf{Нет:} сохраните название исходной функции};
\node[stage,below=27mm of switch](quadrant){Найдите конечную четверть,\ учитывая направление $+\beta$ или $-\beta$};
\node[stage,below=of quadrant](sign){Поставьте знак \textbf{исходной функции} в найденной четверти};
\node[terminus,below=of sign](finish){Подставьте табличное значение и проверьте ОДЗ};
\draw[flow](begin)--(normalize);
\draw[flow](normalize)--(switch);
\draw[flow](switch.south west)--node[left,font=\scriptsize]{да}(yes.north);
\draw[flow](switch.south east)--node[right,font=\scriptsize]{нет}(no.north);
\draw[flow](yes.south)-|([xshift=-1.7cm]quadrant.north);
\draw[flow](no.south)-|([xshift=1.7cm]quadrant.north);
\draw[flow](quadrant)--(sign);
\draw[flow](sign)--(finish);
\end{tikzpicture}
\end{center}
\clearpage

\sectionheading{Тренировочный блок}
\textcolor{muted}{Путь А --- 10 заданий. Сложность повышается постепенно. Указывайте ограничения там, где есть знаменатели.}
\begin{enumerate}[label=\textbf{\arabic*.},itemsep=6pt]
\item Упростите: $4(\sin^2a+\cos^2a)-3$.
\item Известно, что $\sin a=\dfrac8{17}$ и $\dfrac\pi2<a<\pi$.
Найдите $\cos a$, $\tg a$ и $\ctg a$.
\item Вычислите: $\sin15^\circ$.
\item Вычислите: $\cos\dfrac{13\pi}{3}$.
\item Сверните произведение: $\sin4x\cos4x$.
\item Упростите выражение: $\dfrac{\sin6x}{\cos3x}$.
Укажите допустимые значения переменной.
\item Вычислите: $\cos\left(-\dfrac{5\pi}{2}+\dfrac\pi3\right)$.
\item Примените формулу приведения:
$\tg\left(\dfrac\pi2+3a\right)$.
Укажите область определения.
\item Вычислите:
$\sqrt{32}\left(\cos^2\dfrac{5\pi}{8}-\sin^2\dfrac{5\pi}{8}\right)$.
\item Упростите:
$\dfrac{1-\cos2x}{\sin2x}$.
Укажите допустимые значения переменной.
\end{enumerate}
\vspace{4mm}
\textcolor{darkaccent}{\textbf{Правило тренировки.}}
В заданиях 2, 7 и 8 сначала определите четверть или направление поворота;
в заданиях 5, 6, 9 и 10 ищите формулу двойного угла.

\clearpage
\sectionheading{Ответы}
\small
\renewcommand{\arraystretch}{1.43}
\begin{center}
\begin{tabularx}{\linewidth}{@{}>{\bfseries}p{10mm} X@{}}
\toprule
№ & Ответ\\
\midrule
1 & $1$\\
2 & $\cos a=-\dfrac{15}{17}$, $\tg a=-\dfrac8{15}$, $\ctg a=-\dfrac{15}{8}$\\
3 & $\dfrac{\sqrt6-\sqrt2}{4}$\\
4 & $\dfrac12$\\
5 & $\dfrac12\sin8x$\\
6 & $2\sin3x$, при $\cos3x\ne0$\\
7 & $\dfrac{\sqrt3}{2}$\\
8 & $-\ctg3a$, при $\sin3a\ne0$\\
9 & $-4$\\
10 & $\tg x$, при $\sin2x\ne0$ (то есть $x\ne\dfrac{\pi k}{2}$, $k\in\mathbb Z$)\\
\bottomrule
\end{tabularx}
\end{center}
\normalsize
\vspace{8mm}
\sectionheading{Итоговая памятка}
\small
\renewcommand{\arraystretch}{1.35}
\begin{tabularx}{\linewidth}{@{}p{0.36\linewidth}X@{}}
\toprule
Сигнал в условии & Что делать\\
\midrule
Известен синус или косинус & Используйте $\sin^2a+\cos^2a=1$ и выбирайте знак по четверти.\\
Сумма известных углов & Примените формулу суммы или разности.\\
Два одинаковых аргумента & Попробуйте формулу двойного угла в обратном направлении.\\
Угол кратен $\pi/2$ & Проверьте, требуется ли замена на кофункцию.\\
Большой положительный/отрицательный угол & Упростите поворот по модулю $2\pi$, затем определите четверть.\\
Есть знаменатель & До сокращения запишите, когда он отличен от нуля.\\
\bottomrule
\end{tabularx}
\vfill
\begin{center}
\textcolor{muted}{\footnotesize Лёвин Артём Александрович --- эксклюзивно для Дениса Ермолаева}
\end{center}
\end{document}
